\documentclass{chsmc}
\chsmcsetup{
	year = 2024,
	part = 1,
	type = A,
	show-answers = false,
}
\begin{document}

\section{填空题：本大题共 8 小题，每小题 8 分，共 64 分.}

\begin{enumerate}
	\item 若实数 $m > 1$ 满足 $\log_9(\log_8 m) = 2024$，
		则 $\log_3(\log_2 m)$ 的值为 \blankline
	\item 设无穷等比数列 $\{a_n\}$ 的公比 $q$ 满足 $0 < |q| < 1$.
		若 $\{a_n\}$ 的各项和等于 $\{a_n\}$ 各项的平方和，
		则 $a_2$ 的取值范围是 \blankline
	\item 设实数 $a$, $b$ 满足：集合
		$A = \{x\in \mathbf{R} \mid x^2 - 10x + a \leqslant 0\}$ 与
		$B = \{x\in \mathbf{R} \mid bx \leqslant b^3\}$ 的交集为 $[4,9]$，
		则 $a+b$ 的值为 \blankline
	\item 在三棱锥 $P\text{-}ABC$ 中，若 $PA \perp$ 底面 $ABC$，
		且棱 $AB$, $BP$, $BC$, $CP$ 的长分别为 $1$, $2$, $3$, $4$，
		则该三棱锥的体积为 \blankline
	\item 一个不均匀的骰子，掷出 $1$, $2$, $3$, $4$, $5$, $6$ 点的概率依次成等差数列.
		独立地先后掷该骰子两次，所得的点数分别记为 $a$, $b$.
		若事件“$a+b=7$”发生的概率为 $\dfrac{1}{7}$，
		则事件“$a=b$”发生的概率为 \blankline
	\item 设 $f(x)$ 是定义域为 $\mathbf{R}$、最小正周期为 $5$ 的函数.
		若函数 $g(x) = f(2^x)$ 在区间 $[0,5)$ 上的零点个数为 $25$，
		则 $g(x)$ 在区间 $[1,4)$ 上的零点个数为 \blankline
	\item 设 $F_1$, $F_2$ 为椭圆 $\Omega$ 的焦点，在 $\Omega$ 上取一点 $P$
		（异于长轴端点），记 $O$ 为 $\triangle PF_1F_2$ 的外心，若
		$\overrightarrow{PO}\cdot\overrightarrow{F_1F_2} =
		2\overrightarrow{PF_1}\cdot\overrightarrow{PF_2}$，
		则 $\Omega$ 的离心率的最小值为 \blankline
	\item 若三个正整数 $a$, $b$, $c$ 的位数之和为 $8$，且组成 $a$, $b$, $c$
		的 $8$ 个数码能排列为 $2$, $0$, $2$, $4$, $0$, $9$, $0$, $8$，
		则称 $(a,b,c)$ 为“幸运数组”，例如 $(9,8,202400)$ 是一个幸运数组.
		满足 $10 < a < b < c$ 的幸运数组 $(a,b,c)$ 的个数为 \blankline
\end{enumerate}

\section{解答题：本大题共 3 个小题，满分 56 分. 解答应写出文字说明、证明过程或演算步骤.}

\begin{enumerate}[resume]
	\item (本题满分 16 分) 在 $\triangle ABC$ 中，已知
		$\cos C = \dfrac{\sin A + \cos A}{2} = \dfrac{\sin B + \cos B}{2}$，
		求 $\cos C$ 的值.
	\item (本题满分 20 分) 在平面直角坐标系中，双曲线
		$\Gamma\colon x^2-y^2=1$ 的右顶点为 $A$.
		将圆心在 $y$ 轴上，且与 $\Gamma$ 的两支各恰有一个公共点的圆称为“好圆”.
		若两个好圆外切于点 $P$，圆心距为 $d$，求 $\dfrac{d}{|PA|}$ 的所有可能的值.
	\item (本题满分 20 分) 设复数 $z$, $w$ 满足 $z+w=2$，
		求 $S = |z^2-2w| + |w^2-2z|$ 的最小可能值.
\end{enumerate}

\end{document}
